% =============================================================================
% SPECTRA — Section II: Related Work
% Target: IEEE Transactions on Information Forensics & Security
% =============================================================================

\section{Related Work}
\label{sec:related_work}

The SPECTRA framework draws on a broad intellectual heritage spanning graph
theory, deep generative models, density-based clustering, optimal transport,
and Bayesian inference.  This section surveys the most relevant prior art along
each axis, identifies the limitations that motivate our design, and explicitly
positions SPECTRA with respect to the state of the art.

% -----------------------------------------------------------------------------
\subsection{Bitcoin Transaction Graph Analysis}
\label{subsec:bitcoin_graph}
% -----------------------------------------------------------------------------

The Bitcoin blockchain exposes a rich, public ledger of pseudonymous
transactions that has attracted sustained analytical attention since its
inception.  Early graph-theoretic studies established the foundational
properties of this network.  Ron and Shamir~\cite{Ron2013} constructed a
directed graph over approximately 3.3 million Bitcoin addresses and quantified
its degree distribution, finding a power-law tail consistent with scale-free
network topology; they further identified exchange hubs as structurally
dominant nodes whose removal would fragment large connected components.
Meiklejohn et al.~\cite{Meiklejohn2013} introduced the \emph{common-input
ownership heuristic} --- the observation that all addresses co-signing a
transaction likely belong to the same entity --- and demonstrated that
systematic application of this heuristic to address clustering could de-anonymize
a substantial fraction of on-chain flows.  This heuristic remains the
cornerstone of commercial blockchain analytics platforms~\cite{Chainalysis2023}
and provides the entity-resolution substrate on which our own address-level
feature construction rests.

M\"{o}ser et al.~\cite{Moser2017} conducted a large-scale empirical study of
Bitcoin transaction anonymity, characterising privacy-enhancing services and
their effectiveness against graph-based de-anonymisation.  Jourdan et
al.~\cite{Jourdan2019} operationalised a richer feature vocabulary --- combining
topological, temporal, and value-based attributes --- to characterise entity
types in the Bitcoin blockchain, establishing that supervised classification
from local transaction-graph features achieves competitive accuracy on the
Elliptic dataset~\cite{Weber2019}.  Akcora et al.~\cite{Akcora2022} tackled
ransomware prediction through a topological data analysis lens, modeling address
subgraphs as persistent homology descriptors and integrating first-order Markov
chain transition features over address behavioral sequences.

The arrival of large-scale graph neural network (GNN) benchmarks transformed
the field.  Weber et al.~\cite{Weber2019} released the Elliptic dataset --- a
temporal graph of 203,769 Bitcoin transactions with ground-truth illicit/licit
labels --- and demonstrated that a basic graph convolutional network
(GCN)~\cite{Kipf2017} outperforms classical machine learning baselines,
establishing the first large-scale GNN result on real Bitcoin data.  Jha et
al.~\cite{Jha2023} extended this line of work by deploying a temporal graph
network that explicitly models the causal ordering of transaction edges,
reporting improved illicit-transaction recall on the Elliptic benchmark and
confirming the value of dynamic graph representations for forensic
classification; this work appeared in \emph{IEEE Transactions on Information
Forensics and Security}, the same venue we target.  Liu et al.~\cite{Liu2024}
proposed a graph autoencoder architecture for blockchain anomaly detection,
published in \emph{IEEE Transactions on Dependable and Secure Computing},
leveraging reconstruction error in latent space to surface structurally
aberrant subgraphs.  Fan et al.~\cite{Fan2024} introduced AA-GNN, an
anomaly-aware graph neural network for blockchain fraud detection published in
\emph{IEEE Transactions on Industrial Informatics}, augmenting GNN message
passing with an attention mechanism conditioned on predicted anomaly scores to
amplify signals from suspicious neighbourhoods.

\smallskip
\noindent\textbf{Gap addressed by SPECTRA.}
All of the above contributions are conceived as single-task systems: they
perform fraud detection \emph{or} address clustering \emph{or} anomaly scoring,
treating the other objectives as out-of-scope.  None couples spectral graph
embeddings with probabilistic trust estimation and online distribution-shift
monitoring within a single coherent pipeline.  Temporal methods such as
EvolveGCN~\cite{Pareja2020} and BERT4ETH~\cite{He2023} advance the dynamic
modelling of blockchain graphs but do not provide the forensically actionable
multi-output signature that practitioners require.  SPECTRA closes this gap.

% -----------------------------------------------------------------------------
\subsection{Spectral Graph Methods}
\label{subsec:spectral}
% -----------------------------------------------------------------------------

Spectral graph theory provides the theoretical grounding for extracting
geometry-aware node representations from the graph Laplacian.  The seminal
tutorial by Von Luxburg~\cite{VonLuxburg2007} formalises the relationship
between Laplacian eigenvectors and graph connectivity, showing that the $k$
eigenvectors associated with the $k$ smallest non-zero eigenvalues partition the
graph into $k$ weakly-coupled components with minimum normalised cut.  Belkin
and Niyogi~\cite{Belkin2003} demonstrated that Laplacian eigenmaps preserve
neighbourhood geometry under smooth manifold assumptions, establishing a
principled nonlinear dimensionality reduction algorithm that underpins our
spectral embedding module.  In SPECTRA, we compute the top-$d$ eigenvectors of
the symmetrically-normalised Laplacian of the transaction subgraph and use them
as geometric coordinates that supplement transaction-level feature vectors
before ingestion by the VAE encoder.

A practical obstacle to deploying spectral methods on large transaction graphs
is the cost of eigendecomposition: the ARPACK-based \texttt{eigsh} solver
frequently fails to converge on the highly skewed, multi-component graphs that
arise from Bitcoin activity (a bug documented in our codebase).  We resolve this
by adopting the randomised singular value decomposition of Halko, Martinsson,
and Tropp~\cite{Halko2011}, which produces rank-$d$ approximations in
$\mathcal{O}(n d \log d)$ time with probabilistic error guarantees, making
spectral embedding tractable on graphs with millions of nodes.

The emergence of spectral graph convolution networks bridged classical spectral
graph theory with deep learning.  Kipf and Welling~\cite{Kipf2017} simplified
spectral convolution to a first-order Chebyshev approximation, yielding the
widely-used GCN layer $H^{(l+1)} = \sigma\!\left(\tilde{D}^{-\frac{1}{2}}
\tilde{A}\tilde{D}^{-\frac{1}{2}} H^{(l)} W^{(l)}\right)$, where
$\tilde{A} = A + I$ is the adjacency matrix with self-loops.  SPECTRA uses a
two-layer GCN variant (cf.~\cite{Kipf2016b}) within its GraphSAGE-based
classification head.  Hamilton et al.~\cite{Hamilton2017} introduced
GraphSAGE, which learns neighbourhood aggregation functions inductively,
enabling generalisation to nodes unseen during training.  We choose GraphSAGE
over the Graph Attention Network (GAT) of Veli\v{c}kovi\'{c} et
al.~\cite{Velickovic2018} because GAT's per-edge attention computation incurs
$\mathcal{O}(|E|)$ memory overhead that becomes prohibitive on the sparse,
high-degree exchange-hub neighbourhoods characteristic of Bitcoin transaction
graphs; GraphSAGE's mean aggregator achieves comparable accuracy on our dataset
at a fraction of the memory cost.  This design choice is validated empirically
in Section~\ref{sec:experiments}.

% -----------------------------------------------------------------------------
\subsection{Variational Autoencoders and Anomaly Detection}
\label{subsec:vae}
% -----------------------------------------------------------------------------

Variational autoencoders (VAEs), introduced by Kingma and Welling~\cite{Kingma2014},
learn a probabilistic generative model $p_\theta(\mathbf{x}|\mathbf{z})$ jointly
with an approximate posterior $q_\phi(\mathbf{z}|\mathbf{x})$ by maximising the
evidence lower bound (ELBO):
\begin{equation}
  \mathcal{L}_{\text{VAE}} =
      \mathbb{E}_{q_\phi}\!\left[\log p_\theta(\mathbf{x}|\mathbf{z})\right]
      - \beta\,D_{\mathrm{KL}}\!\left(q_\phi(\mathbf{z}|\mathbf{x}) \,\|\,
        p(\mathbf{z})\right),
  \label{eq:elbo}
\end{equation}
where $\beta=1$ recovers the standard VAE and $\beta > 1$ yields the
$\beta$-VAE of Higgins et al.~\cite{Higgins2017}, which encourages
disentangled latent codes.  In SPECTRA we set $\beta = 4$ to promote
interpretable separation of behavioral axes (e.g., volume, fan-out, fee rate)
in the latent space, making the reconstruction score a cleaner proxy for
anomalousness.

The use of VAEs for anomaly detection is well-established.  An and
Cho~\cite{AnCho2015} proposed evaluating anomaly severity via reconstruction
probability --- the probability $p_\theta(\mathbf{x}|\mathbf{z})$ averaged over
samples from $q_\phi(\mathbf{z}|\mathbf{x})$ --- rather than the raw mean squared
error, demonstrating improved performance on the KDD Cup 1999 intrusion
detection dataset.  Xu et al.~\cite{Xu2018} applied a VAE to multivariate KPI
time series in large-scale internet services (DONUT), showing that the
probabilistic reconstruction score is robust to temporal non-stationarity that
defeats threshold-based detectors.  Chalapathy and
Chawla~\cite{Chalapathy2019} survey the broader landscape of deep
anomaly detection, covering autoencoders, GANs, and hybrid
approaches, and conclude that reconstruction-based methods are the most
data-efficient paradigm when labelled anomalies are scarce --- a condition
endemic to blockchain forensics where confirmed illicit labels represent under
two percent of on-chain activity~\cite{Chainalysis2023}.

SPECTRA inherits this reconstruction-probability anomaly score but extends the
VAE's role: the learned latent mean vectors $\boldsymbol{\mu}$ serve
simultaneously as compressed node representations for downstream HDBSCAN
clustering and as inputs to the Bayesian trust scorer, creating a tight coupling
between the generative model and the inference modules described below.

% -----------------------------------------------------------------------------
\subsection{Density-Based Clustering}
\label{subsec:clustering}
% -----------------------------------------------------------------------------

Clustering Bitcoin addresses into behavioural cohorts is a prerequisite for
interpretable forensic reporting.  Classical $k$-means~\cite{Lloyd1982}
requires the analyst to specify $k$ and assumes spherical, equal-variance
clusters, both of which are untenable for the arbitrarily-shaped and variable-density
communities that arise in transaction graphs.  DBSCAN~\cite{Ester1996}
relaxed the sphericity assumption by defining clusters as maximal
density-connected components, but introduced the $\varepsilon$ (neighbourhood
radius) hyperparameter whose sensitivity to scale and density heterogeneity
renders manual tuning impractical on datasets spanning five years of Bitcoin
activity.

HDBSCAN~\cite{McInnes2017}, the hierarchical extension of DBSCAN, eliminates
$\varepsilon$ entirely by constructing a hierarchy of density-connected
components across all scales and extracting the most persistent clusters via the
algorithm's stability criterion.  Campello et al.~\cite{Campello2013} provide
the theoretical foundation for this density-based hierarchical approach,
demonstrating that the resulting partition is optimal in the sense of maximising
the integral of within-cluster density over the hierarchy.  In practice,
HDBSCAN exposes only two sensitive hyperparameters --- \texttt{min\_cluster\_size}
and \texttt{min\_samples} --- both of which have physically interpretable
semantics in our domain (minimum address cohort size and core-point density
threshold respectively).  Crucially, HDBSCAN natively assigns a soft cluster
membership vector per node and identifies noise points (label $-1$), providing
a natural mechanism for flagging structurally isolated addresses that may
warrant human review.

Relative to $k$-means and DBSCAN, HDBSCAN delivers two concrete advantages in
SPECTRA: (i) it discovers the number of behavioural cohorts from data, avoiding
the model-selection overhead of $k$-means; (ii) the noise label serves as an
independent anomaly signal complementary to the VAE reconstruction score.
Together, these properties make HDBSCAN the most appropriate clustering
algorithm for our heterogeneous, large-scale Bitcoin embedding space.

% -----------------------------------------------------------------------------
\subsection{Concept Drift Detection}
\label{subsec:drift}
% -----------------------------------------------------------------------------

Bitcoin transaction patterns are non-stationary: macroeconomic events
(exchange collapses, regulatory announcements), protocol upgrades (SegWit
activation, Taproot deployment), and seasonal trading cycles all induce shifts
in the joint distribution of on-chain features.  A forensic pipeline trained on
historical data must therefore be equipped to detect when its operating
assumptions are violated.

The theoretical framework for quantifying distributional distance used in
SPECTRA is the Wasserstein-1 (Earth Mover's) distance, whose measure-theoretic
foundations are treated comprehensively by Villani~\cite{Villani2009}.  The
Wasserstein distance is particularly appropriate for distributions over discrete
cluster-label histograms because it respects the geometry of the label space
--- two distributions concentrated on adjacent clusters contribute less to the
distance than two distributions concentrated on maximally-separated clusters,
a property that $\chi^2$ divergence and total variation distance do not share.

Survey treatments of concept drift by Lu et al.~\cite{Lu2018} and Gama et
al.~\cite{Gama2014} categorise drift detectors by their statistical assumptions:
window-based methods (Page-Hinkley, ADWIN) monitor univariate statistics and
are sensitive to arbitrary distribution shifts but provide no geometric
information about the nature of the shift; distribution-distance methods
offer richer diagnostics.  Prior blockchain monitoring
systems~\cite{Chainalysis2023} rely on sliding-window anomaly count thresholds
--- a univariate signal blind to compositional changes in the cluster
population.  SPECTRA instead monitors the Wasserstein-1 distance between
consecutive monthly cluster-label distributions, raising a drift flag when this
distance exceeds a control-chart threshold calibrated on the first three months
of the dataset.  This geometry-aware formulation allows us to distinguish, for
example, a benign rotation of activity between exchange and payment cohorts from
a genuine structural shift associated with the emergence of a new transaction
pattern (see Section~\ref{sec:experiments}, drift case studies).

% -----------------------------------------------------------------------------
\subsection{Bayesian Trust and Markov Chain Models}
\label{subsec:bayesian}
% -----------------------------------------------------------------------------

Trust quantification in peer-to-peer and financial networks has a rich Bayesian
pedigree.  J{\o}sang et al.~\cite{Josang2007} formalised subjective logic as
an algebra for uncertainty-aware trust reasoning, enabling the composition of
trust chains in distributed P2P systems while preserving epistemic uncertainty.
The Beta reputation system of J{\o}sang and Ismail~\cite{JosangIsmail2002}
operationalises this framework by maintaining a $\text{Beta}(\alpha, \beta)$
posterior over each agent's trustworthiness, updated via conjugate Bayesian
updating upon each observed transaction outcome; the posterior mean $\alpha /
(\alpha + \beta)$ serves as the trust score and the posterior variance as its
uncertainty.  SPECTRA adopts precisely this conjugate-update scheme for each
Bitcoin address: positive signals (transactions consistent with benign clusters,
low reconstruction error, stable temporal behaviour) increment $\alpha$;
negative signals (high anomaly score, noise cluster assignment, elevated
Wasserstein drift contribution) increment $\beta$.  The resulting per-address
trust score is a calibrated probability with associated credible interval,
making it directly actionable for risk tiering in compliance workflows.

Behavioral sequence modelling via Markov chains has been applied to fraud
detection since the foundational work of Fawcett and Provost~\cite{FawcettProvost1997},
who demonstrated that an address's \emph{sequence} of transaction types (not
merely its aggregate statistics) is a powerful fraud discriminator.  Akcora et
al.~\cite{Akcora2022} applied first-order Markov chain features to Bitcoin
address behavior in the context of ransomware prediction, confirming that
transition probabilities between address-state categories (miner coinbase,
mixing service, exchange, retail) carry forensic signal independent of static
graph features.  SPECTRA constructs per-address Markov transition matrices over
a five-state alphabet derived from cluster assignments and flags addresses whose
stationary distribution diverges significantly from the population-level
stationary distribution --- an information-theoretic test grounded in the
Kullback-Leibler divergence~\cite{Kraskov2004}.

% -----------------------------------------------------------------------------
\subsection{Transformer-Based Sequence Models for Blockchain}
\label{subsec:transformers}
% -----------------------------------------------------------------------------

The success of transformer architectures in natural language processing has
motivated their application to transaction sequence modelling.  He et
al.~\cite{He2023} introduced BERT4ETH, a BERT-style~\cite{Lin2017}
pre-trained transformer for Ethereum account behaviour, treating each account's
transaction history as a sentence of discrete tokens.  Pre-training on the
masked transaction modelling objective and fine-tuning on downstream fraud and
phishing detection tasks, BERT4ETH established state-of-the-art results on
Ethereum-centric benchmarks at WWW 2023.  The concept of treating blockchain
addresses as sequential agents whose transaction histories constitute learnable
token sequences is directly relevant to Bitcoin, and SPECTRA's Markov chain
profiling module can be understood as a lightweight, parameter-free alternative
to the self-attention mechanism that achieves comparable sequence modelling
fidelity without the pre-training cost.

Pareja et al.~\cite{Pareja2020} addressed dynamic graph evolution through
EvolveGCN, which couples an RNN with a GCN to allow model weights to adapt over
time, thus capturing structural non-stationarity without retraining from
scratch.  While EvolveGCN is conceptually aligned with our goal of handling
temporal drift, it treats drift implicitly through weight evolution rather than
explicitly quantifying distributional shift, and it does not produce a
human-interpretable drift signal suitable for regulatory reporting.  SPECTRA's
Wasserstein drift detector provides this explicit, auditable signal.

% -----------------------------------------------------------------------------
\subsection{Positioning of SPECTRA}
\label{subsec:positioning}
% -----------------------------------------------------------------------------

Table~\ref{tab:related_comparison} (Section~\ref{sec:architecture}) summarises
the capabilities of representative prior systems relative to SPECTRA across
five forensically relevant dimensions: spectral embedding, anomaly scoring,
behavioural clustering, drift monitoring, and trust quantification.  No prior
system addresses more than two of these dimensions simultaneously.  The
limitations of existing work can be characterised along three axes:

\smallskip
\noindent\textbf{(1) Task fragmentation.}
Fraud-detection systems~\cite{Weber2019,Jha2023,Fan2024} are optimised for
binary classification and provide no unsupervised clustering or trust scoring.
Anomaly detection systems~\cite{Liu2024,AnCho2015} report scalar anomaly scores
but do not cluster addresses into actionable cohorts or monitor for temporal
distribution shift.  Clustering-centric analyses~\cite{Meiklejohn2013,Jourdan2019}
characterise address populations but offer no continuous trust signal or
automatic drift alert.

\smallskip
\noindent\textbf{(2) Static assumptions.}
Most approaches are trained and evaluated on fixed temporal windows, implicitly
assuming stationarity.  The Bitcoin transaction distribution shifts measurably
across macro-cycles (see Section~\ref{sec:experiments}), and a static model
will degrade silently without a monitoring mechanism.  EvolveGCN~\cite{Pareja2020}
and BERT4ETH~\cite{He2023} partially address non-stationarity but do not expose
a calibrated, auditable drift metric.

\smallskip
\noindent\textbf{(3) Interpretability gap.}
Deep GNN classifiers produce probability scores that are difficult to
audit~\cite{Chalapathy2019}.  Compliance officers require outputs with
associated uncertainty estimates and traceable provenance.  SPECTRA's
Bayesian trust scores carry explicit credible intervals, and each score is
decomposable into its contributing signals (reconstruction error, cluster
assignment, Markov divergence, historical evidence), satisfying the
interpretability requirements of emerging regulatory frameworks.

\smallskip
In contrast to all prior work, SPECTRA is the first framework to tightly couple
spectral graph theory, variational inference, density clustering, Wasserstein
drift monitoring, Bayesian online updating, and Markov chain profiling into a
single differentiable pipeline for Bitcoin transaction analysis, producing four
forensically actionable outputs --- anomaly scores, behavioural cluster
assignments, drift indicators, and calibrated trust scores --- simultaneously,
from the same learned representation.

% =============================================================================
% End of Section II
% =============================================================================
